% =====================================================================
% Ajouts "rapport de recherche" au préambule de la template ECE.
% Chargé par main.tex APRÈS le préambule ECE (donc après babel, csquotes, hyperref).
% Ne pas modifier dans report/V*/ : la source est
% .claude/skills/build-report/assets/preamble-recherche.tex
% =====================================================================

% Typographie et maths
	\usepackage{microtype}
	\usepackage{amsmath,amssymb}
	\usepackage[locale=FR]{siunitx} % \num{0.62} -> 0,62 ; \SI, \qty
	\sisetup{group-minimum-digits=4, separate-uncertainty=true}

% Tableaux
	\usepackage{booktabs} % \toprule \midrule \bottomrule, jamais de traits verticaux
	\usepackage{multirow}
	\usepackage{array}
	\usepackage{tabularx}

% Courts extraits de code (uniquement dans des figures, exemples de perturbation)
	\usepackage{listings}
	\lstset{
		basicstyle=\ttfamily\small,
		columns=fullflexible,
		keepspaces=true,
		breaklines=true,
		frame=single,
		rulecolor=\color{gray!50},
		numbers=none,
		showstringspaces=false,
		extendedchars=true,
		inputencoding=utf8,
		literate={é}{{\'e}}1 {è}{{\`e}}1 {à}{{\`a}}1 {ç}{{\c{c}}}1 {ê}{{\^e}}1 {ù}{{\`u}}1
	}

% Bibliographie (biber)
	\usepackage[backend=biber, style=numeric-comp, sorting=nyt, maxbibnames=6, giveninits=true]{biblatex}
	\addbibresource{refs.bib}

% Références croisées : \cref{fig:x} -> "figure 3". Doit venir après hyperref.
	\usepackage[french,noabbrev]{cleveref}

% Numérotation : les titres affichent \thesection, donc les renvois (\cref) et les titres
% donnent le même numéro (romain pour les parties, lettre pour les annexes).
	\renewcommand{\thesection}{\Roman{section}}
	\renewcommand{\thesubsection}{\thesection.\Alph{subsection}}
	\renewcommand{\thesubsubsection}{\thesubsection.\alph{subsubsection}}
	\titleformat{\section}{\fontsize{18}{18}\bfseries\color{ece}}{\thesection.}{1ex}{}
	\titleformat{\subsection}{\fontsize{16}{16}\bfseries\color{ece}}{\Alph{subsection}.}{1ex}{}
	\titleformat{\subsubsection}{\fontsize{15}{15}\bfseries\color{gamboge}}{(\alph{subsubsection})}{1ex}{}
	\usepackage{appendix}
	\makeatletter
	\g@addto@macro\appendix{%
		\renewcommand{\thesection}{\Alph{section}}%
		\titleformat{\section}{\fontsize{18}{18}\bfseries\color{ece}}{Annexe \thesection.}{1ex}{}%
	}
	\makeatother

% Table des matières
	\setcounter{tocdepth}{2}

% Marqueurs de travail (comptés par `build_report.py check`)
	% \afaire{...} : section ou paragraphe prévu mais sans matière réelle dans le repo.
	\newcommand{\afaire}[1]{%
		\par\vspace{2mm}\noindent
		\fcolorbox{gamboge}{gamboge!10}{\parbox{\dimexpr\linewidth-2\fboxsep-2\fboxrule\relax}{%
			\textbf{\textcolor{gamboge}{À venir.}} #1}}%
		\par\vspace{2mm}}
	% \averifier{...} : affirmation dont la source n'est pas encore vérifiée (ex. chiffre lu sur le web).
	\newcommand{\averifier}[1]{#1\textsuperscript{\textcolor{red}{[à vérifier]}}}

% Résultats chiffrés : jamais tapés à la main.
	% \res{run}{cle} affiche la valeur générée par `build_report.py export-results`
	% dans generated/*.tex. Une clé inconnue s'affiche ?? en rouge et lève un warning.
	\makeatletter
	\newcommand{\res}[2]{%
		\ifcsname res@#1@#2\endcsname
			\csname res@#1@#2\endcsname
		\else
			\textbf{\textcolor{red}{??}}%
			\PackageWarning{codetrace}{Resultat inconnu: #1/#2}%
		\fi}
	\newcommand{\defres}[3]{\expandafter\def\csname res@#1@#2\endcsname{#3}}
	\makeatother
	\IfFileExists{generated/resultats.tex}{\input{generated/resultats.tex}}{}

% Figures et tableaux : légende au-dessus des tableaux, en dessous des figures
	\captionsetup[table]{position=top}
